\documentclass{article}
\usepackage[T1]{fontenc}
\usepackage{lmodern}
\usepackage{graphicx}
\usepackage{amsmath,amssymb}
\usepackage[a4paper,margin=2cm]{geometry}
\usepackage[absolute,overlay]{textpos}
\setlength{\TPHorizModule}{1mm}
\setlength{\TPVertModule}{1mm}
\usepackage[indonesian]{babel}
\usepackage{hyperref}
\setlength{\emergencystretch}{2em}
% unit: Halaman 13

\begin{document}

% === Halaman 13 (python/native) ===

Matriks permutasi kompresi PC-1: 


57 

49 

41 

33 

25 

17 

9 

1 

58 

50 

42 

34 

26 

18 

10 

2 

59 

51 

43 

35 

27 

19 

11 

3 

60 

52 

44 

36 

63 

55 

47 

39 

31 

23 

15 

7 

62 

54 

46 

38 

30 

22 

14 

6 

61 

53 

45 

37 

29 

21 

13 

5 

28 

20 

12 

4 


C0:  berisi bit-bit dari K pada posisi 

57,  49,  41,  33, 25, 17,   9,   1, 58, 50,  42, 34, 26, 18 

10,    2,  59, 51, 43, 35, 27, 19, 11,   3, 60, 52, 44,  36 



D0:  berisi bit-bit dari K pada posisi 

63,  55, 47,  39, 31, 23,  15,   7,  62, 54, 46,  38,  30, 22 

14,    6, 61,  53, 45, 37,  29, 21,  13,   5, 28,  20,  12,   4 

13

Kunci eksternal

Permutasi

PC-1

C0

D0

Left Shift

Left Shift

C1

D1

Left Shift

Left Shift

Permutasi

PC-2

K1

Cj

Dj

Permutasi

PC-2

Kj









Left Shift

Left Shift

C16

D16

Permutasi

PC-2

K16

Perhatikan, bit-bit parity, yaitu bit 

ke-8, 16, 28, … dibuang, sehingga 

menjadi 56 bit

\end{document}
